\documentclass[11pt]{scrartcl}

\usepackage[magyar]{babel}
\usepackage[stixtwo]{fontsetup}
\usepackage{amsmath}
\usepackage{amsthm}
\usepackage{booktabs}
\usepackage[margin=1in]{geometry}
\usepackage{tikz}
\usepackage{pgfplots}

\usetikzlibrary{arrows.meta}
\pgfplotsset{compat=1.18}

% Label sizes. On the website (converted with lwarp) the figures are shown
% scaled down and need larger labels than in the PDF.
\makeatletter
\@ifpackageloaded{lwarp}{%
	\newcommand{\figfont}{\Large}%
	\newcommand{\figtickfont}{\Large}%
	\newcommand{\figdensefont}{\large}%
	\newcommand{\figsmallfont}{\normalsize}%
}{%
	\newcommand{\figfont}{\normalsize}%
	\newcommand{\figtickfont}{\small}%
	\newcommand{\figdensefont}{\small}%
	\newcommand{\figsmallfont}{\scriptsize}%
}
\makeatother

\setlength{\parindent}{0pt}
\setlength{\parskip}{0.66em}

% Egy szín minden függvénynek, az összes ábrán.
\definecolor{sincolor}{HTML}{1F77B4}
\definecolor{coscolor}{HTML}{D62728}
\definecolor{tancolor}{HTML}{2CA02C}
\definecolor{arccolor}{HTML}{FF7F0E}
\definecolor{invcolor}{HTML}{6A3D9A}

\theoremstyle{definition}
\newtheorem{definition}{Definíció}
\newtheorem{theorem}[definition]{Tétel}
\newtheorem{example}[definition]{Példa}
\theoremstyle{remark}
\newtheorem*{remark}{Megjegyzés}

% Közös stílus a függvénygrafikonokhoz.
\pgfplotsset{
	trig/.style={
			axis lines=middle,
			axis line style={-{Stealth[scale=1.1]}, gray!70!black, line width=1pt},
			xlabel={$x$},
			ylabel={$y$},
			xlabel style={anchor=west},
			ylabel style={anchor=south},
			tick style={gray!70!black},
			ticklabel style={font=\figtickfont, fill=white, inner sep=1pt},
			label style={font=\figfont},
			trig format plots=rad,
			% Axes and tick labels (on white boxes) above the curves.
			axis on top,
			samples=200,
			width=\linewidth,
			clip=false,
		},
	pi ticks/.style={
			xtick={-2*pi, -3*pi/2, -pi, -pi/2, pi/2, pi, 3*pi/2, 2*pi},
			xticklabels={$-2\pi$, $-3\pi/2$, $-\pi$, $-\pi/2$,
					$\pi/2$, $\pi$, $3\pi/2$, $2\pi$},
		},
}
\tikzset{
	% Thicker than TikZ's defaults: on the website the figures are shown
	% scaled down, and thin lines get hard to see.
	every picture/.append style={line width=0.9pt, font=\figfont},
	% A curve that leaves the plot gets an arrow tip at that end; one that
	% really ends there gets a dot (end dot).
	tip/.tip={Stealth[length=8pt, width=7pt]},
	continues/.style={tip-tip},
	continues right/.style={-tip},
	end dot/.style={only marks, mark=*, mark size=2.8pt},
	graph/.style={line width=2.2pt, line cap=round},
	faded/.style={line width=1.6pt, gray!55},
	mirror/.style={densely dashed, line width=1pt, gray!75},
}

\title{Trigonometrikus függvények és inverzeik}
\subtitle{Analízis A1 -- segédanyag}
\author{Sándor Tibor}
\date{}

\begin{document}
\maketitle

Ebben a jegyzetben a szinusz, a koszinusz és a tangens függvényt az
egységkörből építjük fel, majd mindhármat invertáljuk. Mivel egyik sem
kölcsönösen egyértelmű, előbb egy alkalmas intervallumra szűkítjük őket; így
kapjuk az arkusz szinusz, az arkusz koszinusz és az arkusz tangens függvényt.

\section{Az egységkör}

Az \emph{egységkör} az origó középpontú, egységnyi sugarú kör:
\[
	x^2 + y^2 = 1\text.
\]
A szögeket \emph{radiánban} mérjük: egy $\alpha$ szög nagysága annak az
ívnek a hossza, amelyet a szög az egységkörből kimetsz. A teljes kör kerülete
$2\pi$, így $360^\circ = 2\pi$, $180^\circ = \pi$ és $90^\circ =
	\frac{\pi}{2}$.

\begin{definition}
	Induljunk el az $(1, 0)$ pontból, és haladjunk az egységkörön $\alpha$
	hosszú utat; ha $\alpha > 0$, az óramutató járásával ellentétesen, ha
	$\alpha < 0$, akkor azzal megegyezően. Jelölje $P_\alpha$ azt a pontot,
	ahová érkezünk. Ekkor
	\[
		\cos\alpha := P_\alpha \text{ első koordinátája,} \qquad
		\sin\alpha := P_\alpha \text{ második koordinátája.}
	\]
	Ha $\cos\alpha \ne 0$, akkor
	\[
		\tan\alpha := \frac{\sin\alpha}{\cos\alpha}\text.
	\]
\end{definition}

A tangensnek is van szemléletes jelentése: az $OP_\alpha$ egyenes az
$x = 1$ érintőt a $(1, \tan\alpha)$ pontban metszi. Ez a hasonló háromszögekből
látszik: a két derékszögű háromszögben a függőleges és a vízszintes befogó
aránya ugyanaz, $\frac{\sin\alpha}{\cos\alpha} = \frac{\tan\alpha}{1}$.

\begin{figure}[ht]
	\centering
	\begin{tikzpicture}[scale=3, >=Stealth, line cap=round]
		\def\ang{40}
		\coordinate (O) at (0, 0);
		\coordinate (P) at (\ang:1);
		\coordinate (C) at ({cos(\ang)}, 0);
		\coordinate (T) at (1, {tan(\ang)});

		\draw[->, gray!70!black] (-1.3, 0) -- (1.5, 0) node[below] {$x$};
		\draw[->, gray!70!black] (0, -1.3) -- (0, 1.3) node[left] {$y$};
		\draw[gray!65] (1, -0.4) -- (1, 1.15) node[above, font=\figtickfont] {$x = 1$};

		\draw[line width=1.5pt] (O) circle (1);
		\draw[tancolor!65] (O) -- (T);
		\draw (O) -- (P) node[midway, above left=-1pt] {$1$};

		\draw[line width=3.4pt, arccolor, opacity=0.8] (1, 0) arc (0:\ang:1);
		\draw[->, arccolor] (0.22, 0) arc (0:\ang:0.22);
		\node[arccolor] at ({\ang / 2}:0.32) {$\alpha$};

		\draw[dashed, gray] (P) -- (0, {sin(\ang)})
		node[left, sincolor] {$\sin\alpha$};
		\draw[line width=2.8pt, coscolor] (O) -- (C)
		node[midway, below] {$\cos\alpha$};
		\draw[line width=2.8pt, sincolor] (C) -- (P);
		\draw[line width=2.8pt, tancolor] (1, 0) -- (T)
		node[midway, right] {$\tan\alpha$};

		\fill (P) circle (0.9pt) node[above=2pt, xshift=-3pt] {$P_\alpha$};
		\fill (T) circle (0.9pt);
		\fill (1, 0) circle (0.9pt) node[below right] {$1$};
		\node[below left] at (O) {$O$};
	\end{tikzpicture}
	\caption{Az $\alpha$ szöghöz tartozó $P_\alpha$ pont az egységkörön. A
		narancssárga ív hossza $\alpha$; a piros szakasz $\cos\alpha$, a kék
		$\sin\alpha$, a zöld $\tan\alpha$.}
	\label{fig:unit-circle}
\end{figure}

\subsection{Alapvető összefüggések}

Mivel $P_\alpha$ rajta van az egységkörön, koordinátái kielégítik a kör
egyenletét. Ez a \emph{trigonometrikus Pitagorasz-tétel}:
\[
	\sin^2\alpha + \cos^2\alpha = 1
	\qquad (\alpha \in \mathbb{R})\text.
\]
Egy teljes kört megtéve ugyanoda érkezünk, ezért $P_{\alpha + 2\pi} =
	P_\alpha$, vagyis
\[
	\sin(\alpha + 2\pi) = \sin\alpha, \qquad
	\cos(\alpha + 2\pi) = \cos\alpha\text.
\]
A többi azonosság a kör szimmetriáiból olvasható le
(\ref{fig:symmetries}.~ábra):
\begin{itemize}
	\item az $x$ tengelyre tükrözve $P_\alpha$-ból $P_{-\alpha}$ lesz:
	      $\cos(-\alpha) = \cos\alpha$ és $\sin(-\alpha) = -\sin\alpha$, azaz a
	      koszinusz páros, a szinusz páratlan;
	\item az $y$ tengelyre tükrözve $P_\alpha$-ból $P_{\pi - \alpha}$ lesz:
	      $\sin(\pi - \alpha) = \sin\alpha$ és $\cos(\pi - \alpha) =
		      -\cos\alpha$;
	\item az origóra tükrözve $P_\alpha$-ból $P_{\alpha + \pi}$ lesz: mindkét
	      koordináta előjelet vált, így $\tan(\alpha + \pi) = \tan\alpha$.
\end{itemize}

\begin{figure}[ht]
	\centering
	\newcommand{\symmetrycircle}[3]{%
		\begin{tikzpicture}[scale=1.35, >=Stealth, line cap=round]
			\def\ang{35}
			\draw[->, gray!70!black] (-1.3, 0) -- (1.35, 0);
			\draw[->, gray!70!black] (0, -1.3) -- (0, 1.3);
			\draw[line width=1.5pt] (0, 0) circle (1);
			\draw[sincolor, line width=1.8pt] (0, 0) -- (\ang:1);
			\draw[coscolor, line width=1.8pt] (0, 0) -- (#1:1);
			\draw[mirror] (\ang:1) -- (#1:1);
			\fill[sincolor] (\ang:1) circle (2pt)
				node[above right=-1pt] {$P_\alpha$};
			\fill[coscolor] (#1:1) circle (2pt) node[#2] {$P_{#3}$};
		\end{tikzpicture}%
	}
	\symmetrycircle{-35}{below right=-1pt}{-\alpha}\hfill
	\symmetrycircle{145}{above left=-1pt}{\pi - \alpha}\hfill
	\symmetrycircle{215}{below left=-1pt}{\alpha + \pi}
	\caption{Tükrözés az $x$ tengelyre, az $y$ tengelyre és az origóra.}
	\label{fig:symmetries}
\end{figure}

\subsection{Nevezetes szögek}

Az alábbi értékeket érdemes fejből tudni. A $\frac{\pi}{4}$ szögnél egy
egyenlő szárú derékszögű háromszög, a $\frac{\pi}{6}$ és a $\frac{\pi}{3}$
szögnél egy félbevágott szabályos háromszög adja az értékeket; a többi
negyedben a szimmetriák segítségével kapjuk meg őket
(\ref{fig:special-angles}.~ábra).

\begin{center}
	\renewcommand{\arraystretch}{2.3}
	\begin{tabular}{@{}l*{5}{c}@{}}
		\toprule
		$\alpha$              & $0$              & $\dfrac{\pi}{6}$      & $\dfrac{\pi}{4}$      &
		$\dfrac{\pi}{3}$      & $\dfrac{\pi}{2}$                                                   \\
		\midrule
		$\sin\alpha$          & $0$              & $\dfrac{1}{2}$        & $\dfrac{\sqrt{2}}{2}$ &
		$\dfrac{\sqrt{3}}{2}$ & $1$                                                                \\
		$\cos\alpha$          & $1$              & $\dfrac{\sqrt{3}}{2}$ & $\dfrac{\sqrt{2}}{2}$ &
		$\dfrac{1}{2}$        & $0$                                                                \\
		$\tan\alpha$          & $0$              & $\dfrac{\sqrt{3}}{3}$ & $1$                   &
		$\sqrt{3}$            & nincs értelmezve                                                   \\
		\bottomrule
	\end{tabular}
\end{center}

\begin{figure}[ht]
	\centering
	\begin{tikzpicture}[scale=2.6, line cap=round, font=\figdensefont]
		\draw[gray!70!black] (-1.08, 0) -- (1.08, 0);
		\draw[gray!70!black] (0, -1.08) -- (0, 1.08);
		\draw[line width=1.5pt] (0, 0) circle (1);
		% Az első negyed pontjainak vetületei.
		\foreach \a in {30, 45, 60} {
				\draw[densely dotted, gray] (\a:1) -- ({cos(\a)}, 0);
				\draw[densely dotted, gray] (\a:1) -- (0, {sin(\a)});
			}
		\foreach \a/\t in {
				0/0, 30/\frac{\pi}{6}, 45/\frac{\pi}{4}, 60/\frac{\pi}{3},
				90/\frac{\pi}{2}, 120/\frac{2\pi}{3}, 135/\frac{3\pi}{4},
				150/\frac{5\pi}{6}, 180/\pi, 210/\frac{7\pi}{6},
				225/\frac{5\pi}{4}, 240/\frac{4\pi}{3}, 270/\frac{3\pi}{2},
				300/\frac{5\pi}{3}, 315/\frac{7\pi}{4}, 330/\frac{11\pi}{6}} {
				\draw[gray!55] (0, 0) -- (\a:1);
				\fill[sincolor] (\a:1) circle (1.3pt);
				\node[fill=white, inner sep=1pt] at (\a:1.25) {$\t$};
			}
		\node[font=\figsmallfont, coscolor, below] at ({sqrt(3) / 2}, 0)
		{$\frac{\sqrt{3}}{2}$};
		\node[font=\figsmallfont, coscolor, below] at (0.5, 0) {$\frac{1}{2}$};
		\node[font=\figsmallfont, sincolor, left] at (0, {sqrt(3) / 2})
		{$\frac{\sqrt{3}}{2}$};
		\node[font=\figsmallfont, sincolor, left] at (0, 0.5) {$\frac{1}{2}$};
	\end{tikzpicture}
	\caption{A nevezetes szögekhez tartozó pontok az egységkörön. Az első
		negyedben a $\frac{\pi}{6}$ és a $\frac{\pi}{3}$ szögű pont vetületei
		is látszanak.}
	\label{fig:special-angles}
\end{figure}

\section{Inverz függvény}

Egy $f\colon A \to B$ függvény
\begin{itemize}
	\item \emph{injektív}, ha különböző pontokban különböző értéket vesz fel:
	      $x_1 \ne x_2 \implies f(x_1) \ne f(x_2)$;
	\item \emph{szürjektív}, ha minden $y \in B$ előáll $f(x)$ alakban, azaz
	      $R_f = B$;
	\item \emph{bijektív} (kölcsönösen egyértelmű), ha injektív és szürjektív
	      is.
\end{itemize}

\begin{definition}
	Ha $f\colon A \to B$ bijektív, akkor minden $y \in B$-hez pontosan egy
	olyan $x \in A$ létezik, amelyre $f(x) = y$. Az így kapott
	$f^{-1}\colon B \to A$, $f^{-1}(y) := x$ függvény $f$ \emph{inverze}. Ekkor
	\[
		f^{-1}\bigl(f(x)\bigr) = x \quad (x \in A)
		\qquad\text{és}\qquad
		f\bigl(f^{-1}(y)\bigr) = y \quad (y \in B)\text.
	\]
\end{definition}

Az inverz grafikonját úgy kapjuk, hogy $f$ grafikonját tükrözzük az $y = x$
egyenesre: ha $(a, b)$ rajta van $f$ grafikonján, akkor $(b, a)$ rajta van
$f^{-1}$ grafikonján.

Egy szigorúan monoton függvény mindig injektív. Ha tehát egy függvény nem
injektív, akkor értelmezési tartományát leszűkíthetjük egy olyan
intervallumra, ahol szigorúan monoton; a képhalmazt pedig az értékkészletre
szűkítve szürjektívvé tehetjük. Pontosan ezt tesszük a trigonometrikus
függvényekkel.

\section{A szinusz és az arkusz szinusz}

\subsection{A szinuszfüggvény}

Ha $P_\alpha$ körbejár az egységkörön, második koordinátája $-1$ és $1$
között hullámzik. Ha ezt a magasságot $\alpha$ függvényében ábrázoljuk, a
szinuszgörbét kapjuk (\ref{fig:sine-from-circle}.~ábra): az egységkör ívét
„kiterítjük” az $x$ tengelyre.

\begin{figure}[ht]
	\centering
	\begin{tikzpicture}[scale=1.25, >=Stealth, line cap=round]
		\def\a{1.0}
		\def\cx{-1.8}
		% Egységkör.
		\draw[->, gray!70!black] (\cx - 1.25, 0) -- (\cx + 1.3, 0);
		\draw[->, gray!70!black] (\cx, -1.3) -- (\cx, 1.35);
		\draw[line width=1.5pt] (\cx, 0) circle (1);
		\draw[line width=3.4pt, arccolor, opacity=0.8]
		(\cx + 1, 0) arc (0:{deg(\a)}:1);
		\draw (\cx, 0) -- ++({deg(\a)}:1);
		\draw[line width=2.8pt, sincolor]
		({\cx + cos(deg(\a))}, 0) -- ++(0, {sin(deg(\a))});
		\fill ({\cx + cos(deg(\a))}, {sin(deg(\a))}) circle (1.9pt)
			node[above=2pt] {$P_\alpha$};
		% Grafikon.
		\draw[->, gray!70!black] (-0.2, 0) -- (6.75, 0) node[below] {$x$};
		\draw[->, gray!70!black] (0, -1.3) -- (0, 1.35) node[left] {$y$};
		\draw[graph, sincolor, line width=2.2pt, continues right]
		plot[domain=0:2*pi, samples=150] (\x, {sin(\x r)});
		\draw[line width=3.4pt, arccolor, opacity=0.8] (0, 0) -- (\a, 0);
		\node[arccolor, below] at ({\a / 2}, 0) {$\alpha$};
		\draw[line width=2.8pt, sincolor] (\a, 0) -- (\a, {sin(\a r)});
		\fill (\a, {sin(\a r)}) circle (1.9pt);
		\draw[mirror] ({\cx + cos(deg(\a))}, {sin(deg(\a))})
		-- (\a, {sin(\a r)});
		\foreach \x/\t in {
				1.5708/{\pi/2}, 3.1416/\pi, 4.7124/{3\pi/2}, 6.2832/2\pi} {
				\draw[gray!70!black] ({\x}, 0.05) -- ({\x}, -0.05)
					node[below, font=\figtickfont] {$\t$};
			}
		\draw[gray!70!black] (0.05, 1) -- (-0.05, 1)
		node[left, font=\figtickfont] {$1$};
		\draw[gray!70!black] (0.05, -1) -- (-0.05, -1)
		node[left, font=\figtickfont] {$-1$};
	\end{tikzpicture}
	\caption{A szinuszgörbe az egységkörből: a $P_\alpha$ pont magassága az
		$x = \alpha$ helyen.}
	\label{fig:sine-from-circle}
\end{figure}

Az egységkörös definícióból és az előző fejezet azonosságaiból a szinusz
főbb tulajdonságai:
\begin{itemize}
	\item $D_{\sin} = \mathbb{R}$ és $R_{\sin} = [-1, 1]$;
	\item $2\pi$ szerint periodikus: $\sin(x + 2\pi) = \sin x$;
	\item páratlan: $\sin(-x) = -\sin x$, grafikonja az origóra szimmetrikus;
	\item zérushelyei: $\sin x = 0 \iff x = k\pi$ $(k \in \mathbb{Z})$.
\end{itemize}

\begin{figure}[ht]
	\centering
	\begin{tikzpicture}
		\begin{axis}[
				trig, pi ticks,
				height=5.5cm,
				xmin=-2*pi - 0.4, xmax=2*pi + 1,
				ymin=-1.4, ymax=1.5,
				ytick={-1, 1},
			]
			\addplot[graph, continues, sincolor, domain=-2*pi - 0.3:2*pi + 0.3] {sin(x)}
			node[pos=0.6, above right] {$\sin x$};
			\addplot[mirror, continues, domain=-2*pi - 0.3:2*pi + 0.3] {0.5}
			node[right, black] {$y = 1/2$};
			\foreach \x in {-11*pi/6, -7*pi/6, pi/6, 5*pi/6} {
					\edef\temp{\noexpand\addplot[only marks, mark=*,
							mark size=2.8pt, black] coordinates {(\x, 0.5)};}
					\temp
				}
		\end{axis}
	\end{tikzpicture}
	\caption{A szinusz nem injektív: az $y = \frac{1}{2}$ egyenes végtelen
		sok pontban metszi a grafikonját.}
	\label{fig:sine-not-injective}
\end{figure}

\subsection{Bijektívvé tétel}

A szinusz nem injektív, hiszen például
\[
	\sin\frac{\pi}{6} = \sin\frac{5\pi}{6} = \sin\frac{13\pi}{6} = \frac{1}{2}
\]
(\ref{fig:sine-not-injective}.~ábra), így nem is invertálható. A
$\left[-\frac{\pi}{2}, \frac{\pi}{2}\right]$ intervallumon viszont szigorúan
monoton nő, és közben minden $[-1, 1]$-beli értéket pontosan egyszer vesz fel
(\ref{fig:sine-restriction}.~ábra). Ezért a
\[
	\sin\colon \left[-\frac{\pi}{2}, \frac{\pi}{2}\right] \to [-1, 1]
\]
leszűkítés bijektív. Más intervallumot is választhatnánk (például
$\left[\frac{\pi}{2}, \frac{3\pi}{2}\right]$-t), de ez a megszokott választás:
ez tartalmazza a $0$-t, és ezen a szinusz növekvő.

\begin{figure}[ht]
	\centering
	\begin{tikzpicture}
		\begin{axis}[
				trig, pi ticks,
				height=5.5cm,
				xmin=-2*pi - 0.4, xmax=2*pi + 1,
				ymin=-1.4, ymax=1.5,
				ytick={-1, 1},
			]
			\fill[sincolor, opacity=0.08] (-pi/2, -1.4) rectangle (pi/2, 1.4);
			\addplot[faded, continues, domain=-2*pi - 0.3:2*pi + 0.3] {sin(x)};
			\addplot[graph, sincolor, line width=2.8pt, domain=-pi/2:pi/2] {sin(x)};
			\addplot[only marks, mark=*, mark size=2.8pt, sincolor]
			coordinates {(-pi/2, -1) (pi/2, 1)};
		\end{axis}
	\end{tikzpicture}
	\caption{A szinusz leszűkítése a $\left[-\frac{\pi}{2},
				\frac{\pi}{2}\right]$ intervallumra már szigorúan monoton.}
	\label{fig:sine-restriction}
\end{figure}

\subsection{Az arkusz szinusz}

\begin{definition}
	Az \emph{arkusz szinusz} a leszűkített szinusz inverze:
	\[
		\arcsin\colon [-1, 1] \to \left[-\frac{\pi}{2}, \frac{\pi}{2}\right],
		\qquad
		\arcsin y = x \iff \sin x = y \text{ és }
		x \in \left[-\frac{\pi}{2}, \frac{\pi}{2}\right]\text.
	\]
\end{definition}

Szavakban: $\arcsin y$ az a $\left[-\frac{\pi}{2}, \frac{\pi}{2}\right]$-beli
szög, amelynek szinusza $y$. Grafikonját a leszűkített szinusz grafikonjának
az $y = x$ egyenesre vett tükörképe adja (\ref{fig:arcsin}.~ábra). Az inverz
tulajdonságai a szinuszéból következnek:
\begin{itemize}
	\item $D_{\arcsin} = [-1, 1]$ és
	      $R_{\arcsin} = \left[-\frac{\pi}{2}, \frac{\pi}{2}\right]$;
	\item szigorúan monoton nő;
	\item páratlan: $\arcsin(-y) = -\arcsin y$.
\end{itemize}

\begin{figure}[ht]
	\centering
	\begin{tikzpicture}
		\begin{axis}[
				trig,
				width=0.78\linewidth,
				axis equal image,
				xmin=-2.1, xmax=2.1, ymin=-2.1, ymax=2.1,
				xtick={-pi/2, -1, 1, pi/2},
				xticklabels={$-\pi/2$, $-1$, $1$, $\pi/2$},
				ytick={-pi/2, -1, 1, pi/2},
				yticklabels={$-\pi/2$, $-1$, $1$, $\pi/2$},
			]
			\addplot[mirror, continues, domain=-2:2] {x} node[pos=0.96, below right]
			{$y = x$};
			\addplot[graph, sincolor!75, domain=-pi/2:pi/2] {sin(x)}
			node[pos=0.98, below right, sincolor] {$\sin x$};
			\addplot[end dot, sincolor!75] coordinates {(-pi/2, -1) (pi/2, 1)};
			\addplot[end dot, invcolor] coordinates {(-1, -pi/2) (1, pi/2)};
			\addplot[graph, invcolor, domain=-pi/2:pi/2] ({sin(x)}, {x})
			node[pos=0.88, left, invcolor] {$\arcsin x$};
			\addplot[densely dotted, line width=1.3pt, gray!80!black] coordinates
				{(1, {sin(1)}) ({sin(1)}, 1)};
			\addplot[only marks, mark=*, mark size=2.6pt, sincolor!75]
			coordinates {(1, {sin(1)})};
			\addplot[only marks, mark=*, mark size=2.6pt, invcolor]
			coordinates {({sin(1)}, 1)};
		\end{axis}
	\end{tikzpicture}
	\caption{Az arkusz szinusz grafikonja (lila) a leszűkített szinusz
		grafikonjának (halványkék) tükörképe az $y = x$ egyenesre.}
	\label{fig:arcsin}
\end{figure}

\begin{remark}
	A $\sin(\arcsin y) = y$ egyenlőség minden $y \in [-1, 1]$ esetén igaz, az
	$\arcsin(\sin x) = x$ viszont \emph{csak}
	$x \in \left[-\frac{\pi}{2}, \frac{\pi}{2}\right]$ esetén! Az arkusz
	szinusz ugyanis mindig $\left[-\frac{\pi}{2}, \frac{\pi}{2}\right]$-beli
	értéket ad.
\end{remark}

\begin{example}
	Számítsuk ki az $\arcsin\left(\sin\frac{5\pi}{6}\right)$ értéket!

	Mivel $\frac{5\pi}{6} \notin \left[-\frac{\pi}{2}, \frac{\pi}{2}\right]$,
	a válasz nem $\frac{5\pi}{6}$. Számoljunk lépésenként:
	\[
		\arcsin\left(\sin\frac{5\pi}{6}\right)
		= \arcsin\frac{1}{2}
		= \frac{\pi}{6}\text,
	\]
	hiszen $\sin\frac{\pi}{6} = \frac{1}{2}$ és
	$\frac{\pi}{6} \in \left[-\frac{\pi}{2}, \frac{\pi}{2}\right]$.
\end{example}

\begin{example}
	Mutassuk meg, hogy $\cos(\arcsin y) = \sqrt{1 - y^2}$ minden
	$y \in [-1, 1]$ esetén!

	Legyen $x = \arcsin y$, ekkor $\sin x = y$ és
	$x \in \left[-\frac{\pi}{2}, \frac{\pi}{2}\right]$. Ezen az intervallumon a
	koszinusz nemnegatív, így a trigonometrikus Pitagorasz-tételből
	\[
		\cos(\arcsin y)
		= \cos x
		= \sqrt{1 - \sin^2 x}
		= \sqrt{1 - y^2}\text.
	\]
\end{example}

\section{A koszinusz és az arkusz koszinusz}

\subsection{A koszinuszfüggvény}

A koszinusz $P_\alpha$ első koordinátája. Ahogy a szinuszgörbét a $P_\alpha$
pont magasságából kaptuk, úgy a koszinusz grafikonja az első koordinátájából
adódik: most az egységkör ívét lefelé terítjük ki, és minden $\alpha$ helyen
vízszintesen mérjük fel $\cos\alpha$-t (\ref{fig:cosine-from-circle}.~ábra). Így
a koszinuszgörbét elforgatva látjuk: a függőleges tengely mentén lefelé nő
$\alpha$, és a görbe vízszintesen hullámzik $-1$ és $1$ között.

\begin{figure}[ht]
	\centering
	\begin{tikzpicture}[scale=1.15, >=Stealth, line cap=round]
		\def\a{1.0}
		\def\gap{1.8}
		% Unit circle.
		\draw[->, gray!70!black] (-1.25, 0) -- (1.3, 0);
		\draw[->, gray!70!black] (0, -1.3) -- (0, 1.35);
		\draw[line width=1.5pt] (0, 0) circle (1);
		\draw[line width=3.4pt, arccolor, opacity=0.8]
		(1, 0) arc (0:{deg(\a)}:1);
		\draw (0, 0) -- ({deg(\a)}:1);
		\draw[line width=2.8pt, coscolor] (0, 0) -- ({cos(deg(\a))}, 0);
		\draw[line width=2.8pt, sincolor]
		({cos(deg(\a))}, 0) -- ({deg(\a)}:1);
		\fill ({deg(\a)}:1) circle (1.9pt) node[above=2pt] {$P_\alpha$};

		% Sine, unrolled to the right: the height of P_alpha.
		\begin{scope}[shift={(\gap, 0)}]
			\draw[->, gray!70!black] (-0.2, 0) -- (6.75, 0) node[below] {$x$};
			\draw[->, gray!70!black] (0, -1.3) -- (0, 1.35) node[left] {$y$};
			\draw[graph, sincolor, continues right]
			plot[domain=0:2*pi, samples=150] (\x, {sin(\x r)});
			\draw[line width=3.4pt, arccolor, opacity=0.8] (0, 0) -- (\a, 0);
			\node[arccolor, below] at ({\a / 2}, 0) {$\alpha$};
			\draw[line width=2.8pt, sincolor] (\a, 0) -- (\a, {sin(\a r)});
			\fill (\a, {sin(\a r)}) circle (1.9pt);
			\foreach \x/\t in {
					1.5708/{\pi/2}, 3.1416/\pi, 4.7124/{3\pi/2}, 6.2832/2\pi} {
					\draw[gray!70!black] ({\x}, 0.05) -- ({\x}, -0.05)
					node[below, font=\figtickfont] {$\t$};
				}
			\draw[gray!70!black] (0.05, 1) -- (-0.05, 1)
			node[left, font=\figtickfont] {$1$};
			\draw[gray!70!black] (0.05, -1) -- (-0.05, -1)
			node[left, font=\figtickfont] {$-1$};
		\end{scope}
		\draw[mirror] ({deg(\a)}:1) -- ({\gap + \a}, {sin(\a r)});

		% Cosine, unrolled downwards: the first coordinate of P_alpha.
		\begin{scope}[shift={(0, -\gap)}]
			\draw[->, gray!70!black] (-1.3, 0) -- (1.4, 0) node[right] {$y$};
			\draw[->, gray!70!black] (0, 0.2) -- (0, -6.75) node[left] {$x$};
			\draw[graph, coscolor, continues right]
			plot[domain=0:2*pi, samples=150] ({cos(\x r)}, -\x);
			\draw[line width=3.4pt, arccolor, opacity=0.8] (0, 0) -- (0, -\a);
			\node[arccolor, left] at (0, {-\a / 2}) {$\alpha$};
			\draw[line width=2.8pt, coscolor] (0, -\a) -- ({cos(\a r)}, -\a);
			\fill ({cos(\a r)}, -\a) circle (1.9pt);
			\foreach \y/\t in {
					1.5708/{\pi/2}, 3.1416/\pi, 4.7124/{3\pi/2}, 6.2832/2\pi} {
					\draw[gray!70!black] (0.05, -\y) -- (-0.05, -\y)
					node[left, font=\figtickfont, fill=white, inner sep=1pt]
						{$\t$};
				}
			\draw[gray!70!black] (1, -0.05) -- (1, 0.05)
			node[above, font=\figtickfont] {$1$};
			\draw[gray!70!black] (-1, -0.05) -- (-1, 0.05)
			node[above, font=\figtickfont] {$-1$};
		\end{scope}
		\draw[mirror] ({cos(deg(\a))}, 0) -- ({cos(\a r)}, {-\gap - \a});
	\end{tikzpicture}
	\caption{A szinusz- és a koszinuszgörbe az egységkörből. A $P_\alpha$ pont
		magassága (kék) jobbra, első koordinátája (piros) lefelé kiterítve
		adja a két grafikont; a narancssárga szakaszok hossza $\alpha$.}
	\label{fig:cosine-from-circle}
\end{figure}

Az egységkörön negyed fordulattal elforgatva azt is kapjuk, hogy
\[
	\cos x = \sin\left(x + \frac{\pi}{2}\right)\text,
\]
vagyis a koszinusz grafikonja a szinuszgörbe $\frac{\pi}{2}$-vel balra
eltolva (\ref{fig:cosine}.~ábra). Főbb tulajdonságai:
\begin{itemize}
	\item $D_{\cos} = \mathbb{R}$ és $R_{\cos} = [-1, 1]$;
	\item $2\pi$ szerint periodikus: $\cos(x + 2\pi) = \cos x$;
	\item páros: $\cos(-x) = \cos x$, grafikonja az $y$ tengelyre
	      szimmetrikus;
	\item zérushelyei: $\cos x = 0 \iff x = \frac{\pi}{2} + k\pi$
	      $(k \in \mathbb{Z})$.
\end{itemize}

\begin{figure}[ht]
	\centering
	\begin{tikzpicture}
		\begin{axis}[
				trig, pi ticks,
				height=5.5cm,
				xmin=-2*pi - 0.4, xmax=2*pi + 1,
				ymin=-1.4, ymax=1.5,
				ytick={-1, 1},
			]
			\addplot[faded, continues, sincolor!50, domain=-2*pi - 0.3:2*pi + 0.3] {sin(x)};
			\addplot[graph, continues, coscolor, domain=-2*pi - 0.3:2*pi + 0.3] {cos(x)}
			node[pos=0.02, above right] {$\cos x$};
			\addplot[mirror, continues, domain=-2*pi - 0.3:2*pi + 0.3] {0.5}
			node[right, black] {$y = 1/2$};
			\foreach \x in {-5*pi/3, -pi/3, pi/3, 5*pi/3} {
					\edef\temp{\noexpand\addplot[only marks, mark=*,
							mark size=2.8pt, black] coordinates {(\x, 0.5)};}
					\temp
				}
		\end{axis}
	\end{tikzpicture}
	\caption{A koszinusz grafikonja (piros) a szinuszgörbe (halványkék)
		$\frac{\pi}{2}$-vel balra tolva. A koszinusz sem injektív.}
	\label{fig:cosine}
\end{figure}

\subsection{Bijektívvé tétel}

A koszinusz sem injektív: például
$\cos\left(-\frac{\pi}{3}\right) = \cos\frac{\pi}{3} = \frac{1}{2}$. Mivel
páros, a $0$-t tartalmazó szimmetrikus intervallum most nem jó. A
$[0, \pi]$ intervallumon viszont szigorúan monoton \emph{csökken}, és minden
$[-1, 1]$-beli értéket pontosan egyszer vesz fel (\ref{fig:cosine-restriction}.~ábra).
Ezért a
\[
	\cos\colon [0, \pi] \to [-1, 1]
\]
leszűkítés bijektív.

\begin{figure}[ht]
	\centering
	\begin{tikzpicture}
		\begin{axis}[
				trig, pi ticks,
				height=5.5cm,
				xmin=-2*pi - 0.4, xmax=2*pi + 1,
				ymin=-1.4, ymax=1.5,
				ytick={-1, 1},
			]
			\fill[coscolor, opacity=0.08] (0, -1.4) rectangle (pi, 1.4);
			\addplot[faded, continues, domain=-2*pi - 0.3:2*pi + 0.3] {cos(x)};
			\addplot[graph, coscolor, line width=2.8pt, domain=0:pi] {cos(x)};
			\addplot[only marks, mark=*, mark size=2.8pt, coscolor]
			coordinates {(0, 1) (pi, -1)};
		\end{axis}
	\end{tikzpicture}
	\caption{A koszinusz leszűkítése a $[0, \pi]$ intervallumra már szigorúan
		monoton.}
	\label{fig:cosine-restriction}
\end{figure}

\subsection{Az arkusz koszinusz}

\begin{definition}
	Az \emph{arkusz koszinusz} a leszűkített koszinusz inverze:
	\[
		\arccos\colon [-1, 1] \to [0, \pi],
		\qquad
		\arccos y = x \iff \cos x = y \text{ és } x \in [0, \pi]\text.
	\]
\end{definition}

Tulajdonságai (\ref{fig:arccos}.~ábra):
\begin{itemize}
	\item $D_{\arccos} = [-1, 1]$ és $R_{\arccos} = [0, \pi]$;
	\item szigorúan monoton csökken;
	\item nem páros és nem páratlan, de grafikonja a
	      $\left(0, \frac{\pi}{2}\right)$ pontra szimmetrikus:
	      $\arccos(-y) = \pi - \arccos y$.
\end{itemize}

\begin{figure}[ht]
	\centering
	\begin{tikzpicture}
		\begin{axis}[
				trig,
				width=0.78\linewidth,
				axis equal image,
				xmin=-1.6, xmax=3.6, ymin=-1.6, ymax=3.6,
				xtick={-1, 1, pi/2, pi},
				xticklabels={$-1$, $1$, $\pi/2$, $\pi$},
				ytick={-1, 1, pi/2, pi},
				yticklabels={$-1$, $1$, $\pi/2$, $\pi$},
			]
			\addplot[mirror, continues, domain=-1.5:3.5] {x} node[pos=0.97, below right]
			{$y = x$};
			\addplot[graph, coscolor!75, domain=0:pi] {cos(x)}
			node[pos=1, below right, coscolor] {$\cos x$};
			\addplot[end dot, coscolor!75] coordinates {(0, 1) (pi, -1)};
			\addplot[end dot, invcolor] coordinates {(1, 0) (-1, pi)};
			\addplot[graph, invcolor, domain=0:pi] ({cos(x)}, {x})
			node[pos=1, left, invcolor] {$\arccos x$};
			\addplot[densely dotted, line width=1.3pt, gray!80!black] coordinates
				{(1, {cos(1)}) ({cos(1)}, 1)};
			\addplot[only marks, mark=*, mark size=2.6pt, coscolor!75]
			coordinates {(1, {cos(1)})};
			\addplot[only marks, mark=*, mark size=2.6pt, invcolor]
			coordinates {({cos(1)}, 1)};
		\end{axis}
	\end{tikzpicture}
	\caption{Az arkusz koszinusz grafikonja (lila) a leszűkített koszinusz
		grafikonjának (halványpiros) tükörképe az $y = x$ egyenesre.}
	\label{fig:arccos}
\end{figure}

\begin{theorem}
	Minden $y \in [-1, 1]$ esetén
	\[
		\arcsin y + \arccos y = \frac{\pi}{2}\text.
	\]
\end{theorem}

\begin{proof}
	Legyen $x = \arcsin y \in \left[-\frac{\pi}{2}, \frac{\pi}{2}\right]$.
	Ekkor $\frac{\pi}{2} - x \in [0, \pi]$, és
	\[
		\cos\left(\frac{\pi}{2} - x\right) = \sin x = y\text,
	\]
	tehát $\arccos y = \frac{\pi}{2} - x = \frac{\pi}{2} - \arcsin y$.
\end{proof}

\begin{example}
	Számítsuk ki az $\arccos\left(\cos\left(-\frac{\pi}{3}\right)\right)$
	értéket!

	A $-\frac{\pi}{3}$ nincs a $[0, \pi]$ intervallumban, ezért ismét
	lépésenként számolunk:
	\[
		\arccos\left(\cos\left(-\frac{\pi}{3}\right)\right)
		= \arccos\frac{1}{2}
		= \frac{\pi}{3}\text.
	\]
\end{example}

\section{A tangens és az arkusz tangens}

\subsection{A tangensfüggvény}

A tangens ott van értelmezve, ahol a koszinusz nem nulla:
\[
	\tan x = \frac{\sin x}{\cos x}, \qquad
	D_{\tan} = \left\{ x \in \mathbb{R} : x \ne \frac{\pi}{2} + k\pi,\
	k \in \mathbb{Z} \right\}\text.
\]
Az egységkörön ez az $OP_\alpha$ egyenes és az $x = 1$ érintő metszéspontjának
magassága (\ref{fig:unit-circle}.~ábra). Ha $\alpha \to \frac{\pi}{2}$ alulról,
az egyenes egyre meredekebb, és a metszéspont a végtelenbe tart; ezért a
tangens grafikonjának függőleges aszimptotái vannak az
$x = \frac{\pi}{2} + k\pi$ helyeken (\ref{fig:tangent}.~ábra). Főbb
tulajdonságai:
\begin{itemize}
	\item $R_{\tan} = \mathbb{R}$;
	\item $\pi$ szerint periodikus: $\tan(x + \pi) = \tan x$;
	\item páratlan: $\tan(-x) = -\tan x$;
	\item zérushelyei: $\tan x = 0 \iff x = k\pi$ $(k \in \mathbb{Z})$.
\end{itemize}

\begin{figure}[ht]
	\centering
	\begin{tikzpicture}
		\begin{axis}[
				trig, pi ticks,
				height=7cm,
				xmin=-3*pi/2 - 0.5, xmax=3*pi/2 + 0.6,
				ymin=-4.2, ymax=4.2,
				ytick=\empty,
				restrict y to domain=-4.2:4.2,
			]
			\fill[tancolor, opacity=0.08] (-pi/2, -4.2) rectangle (pi/2, 4.2);
			\foreach \k in {-3, -1, 1, 3} {
					\edef\temp{\noexpand\draw[mirror]
						({\k*pi/2}, -4.2) -- ({\k*pi/2}, 4.2);}
					\temp
				}
			\addplot[faded, continues, domain=-pi - 1.3365:-pi + 1.3365] {tan(x)};
			\addplot[faded, continues, domain=pi - 1.3365:pi + 1.3365] {tan(x)};
			\addplot[graph, continues, tancolor, line width=2.8pt,
				domain=-1.3365:1.3365] {tan(x)}
			node[pos=0.93, right] {$\tan x$};
		\end{axis}
	\end{tikzpicture}
	\caption{A tangens grafikonja. A szaggatott függőleges egyenesek az
		aszimptoták; a $\left(-\frac{\pi}{2}, \frac{\pi}{2}\right)$
		intervallumon (zöld) a tangens szigorúan monoton nő.}
	\label{fig:tangent}
\end{figure}

\subsection{Bijektívvé tétel}

A tangens $\pi$ szerint periodikus, így nem injektív. A
$\left(-\frac{\pi}{2}, \frac{\pi}{2}\right)$ \emph{nyílt} intervallumon
viszont szigorúan monoton nő, és minden valós értéket pontosan egyszer vesz
fel. A végpontokat nem vehetjük hozzá, mert ott a tangens nincs értelmezve.
Így a
\[
	\tan\colon \left(-\frac{\pi}{2}, \frac{\pi}{2}\right) \to \mathbb{R}
\]
leszűkítés bijektív.

\subsection{Az arkusz tangens}

\begin{definition}
	Az \emph{arkusz tangens} a leszűkített tangens inverze:
	\[
		\arctan\colon \mathbb{R} \to \left(-\frac{\pi}{2}, \frac{\pi}{2}\right),
		\qquad
		\arctan y = x \iff \tan x = y \text{ és }
		x \in \left(-\frac{\pi}{2}, \frac{\pi}{2}\right)\text.
	\]
\end{definition}

A tükrözés során a tangens függőleges aszimptotáiból vízszintes aszimptoták
lesznek (\ref{fig:arctan}.~ábra). Tulajdonságai:
\begin{itemize}
	\item $D_{\arctan} = \mathbb{R}$ és
	      $R_{\arctan} = \left(-\frac{\pi}{2}, \frac{\pi}{2}\right)$;
	\item szigorúan monoton nő és páratlan;
	\item $\displaystyle\lim_{y \to \infty} \arctan y = \frac{\pi}{2}$ és
	      $\displaystyle\lim_{y \to -\infty} \arctan y = -\frac{\pi}{2}$.
\end{itemize}

\begin{figure}[ht]
	\centering
	\begin{tikzpicture}
		\begin{axis}[
				trig,
				width=0.75\linewidth,
				axis equal image,
				xmin=-4.6, xmax=4.6, ymin=-4.2, ymax=4.2,
				xtick={-pi/2, pi/2},
				xticklabels={$-\pi/2$, $\pi/2$},
				ytick={-pi/2, pi/2},
				yticklabels={$-\pi/2$, $\pi/2$},
				restrict y to domain=-4.2:4.2,
			]
			\addplot[mirror, continues, domain=-4.2:4.2] {x} node[pos=0.97, below right]
			{$y = x$};
			\draw[mirror, invcolor!60] (-4.5, -pi/2) -- (4.5, -pi/2);
			\draw[mirror, invcolor!60] (-4.5, pi/2) -- (4.5, pi/2);
			\draw[mirror, tancolor!75] (-pi/2, -4.2) -- (-pi/2, 4.2);
			\draw[mirror, tancolor!75] (pi/2, -4.2) -- (pi/2, 4.2);
			\addplot[graph, continues, tancolor!75, domain=-1.3365:1.3365] {tan(x)}
			node[pos=1, right, tancolor] {$\tan x$};
			\addplot[graph, continues, invcolor, domain=-4.5:4.5] {atan(x)}
			node[pos=0.97, below left, invcolor] {$\arctan x$};
		\end{axis}
	\end{tikzpicture}
	\caption{Az arkusz tangens grafikonja (lila) a leszűkített tangens
		grafikonjának (halványzöld) tükörképe. A tangens függőleges
		aszimptotáiból az arkusz tangens vízszintes aszimptotái lesznek.}
	\label{fig:arctan}
\end{figure}

\begin{example}
	Számítsuk ki az $\arctan\left(\tan\frac{3\pi}{4}\right)$ értéket!

	A tangens $\pi$ szerint periodikus, és
	$\frac{3\pi}{4} - \pi = -\frac{\pi}{4} \in \left(-\frac{\pi}{2},
		\frac{\pi}{2}\right)$, így
	\[
		\arctan\left(\tan\frac{3\pi}{4}\right)
		= \arctan\left(\tan\left(-\frac{\pi}{4}\right)\right)
		= -\frac{\pi}{4}\text.
	\]
\end{example}

\begin{example}
	Mutassuk meg, hogy $\cos(\arctan y) = \dfrac{1}{\sqrt{1 + y^2}}$ minden
	$y \in \mathbb{R}$ esetén!

	Legyen $x = \arctan y$, ekkor $\tan x = y$ és
	$x \in \left(-\frac{\pi}{2}, \frac{\pi}{2}\right)$, ahol $\cos x > 0$. A
	trigonometrikus Pitagorasz-tételt $\cos^2 x$-szel elosztva
	\[
		\tan^2 x + 1 = \frac{1}{\cos^2 x}\text,
	\]
	amiből
	\[
		\cos(\arctan y) = \cos x = \frac{1}{\sqrt{1 + \tan^2 x}}
		= \frac{1}{\sqrt{1 + y^2}}\text.
	\]
\end{example}

\section{Összefoglalás}

\begin{center}
	\renewcommand{\arraystretch}{1.7}
	\begin{tabular}{@{}cc@{}}
		\toprule
		Bijektív leszűkítés                                                    &
		Inverz                                                                    \\
		\midrule
		$\sin\colon \left[-\frac{\pi}{2}, \frac{\pi}{2}\right] \to [-1, 1]$    &
		$\arcsin\colon [-1, 1] \to \left[-\frac{\pi}{2}, \frac{\pi}{2}\right]$    \\
		$\cos\colon [0, \pi] \to [-1, 1]$                                      &
		$\arccos\colon [-1, 1] \to [0, \pi]$                                      \\
		$\tan\colon \left(-\frac{\pi}{2}, \frac{\pi}{2}\right) \to \mathbb{R}$ &
		$\arctan\colon \mathbb{R} \to \left(-\frac{\pi}{2}, \frac{\pi}{2}\right)$ \\
		\bottomrule
	\end{tabular}
\end{center}

Mindhárom esetben ugyanazt a receptet követtük:
\begin{enumerate}
	\item megkerestünk egy intervallumot, ahol a függvény szigorúan monoton
	      (így injektív), és minden értékét felveszi;
	\item a képhalmazt az értékkészletre szűkítettük (így szürjektív);
	\item a kapott bijekciót invertáltuk; az inverz grafikonja a leszűkített
	      függvény grafikonjának tükörképe az $y = x$ egyenesre.
\end{enumerate}
Az inverz függvénybe helyettesítés mindig visszaadja az eredeti értéket
($\sin(\arcsin y) = y$), fordítva azonban csak a leszűkített intervallumon
($\arcsin(\sin x) = x$ csak ha $x \in \left[-\frac{\pi}{2},
		\frac{\pi}{2}\right]$).
\end{document}
